problem:  
Let \( \Sigma^n \) (\( n \ge 7 \)) be an exotic sphere—i.e., a smooth manifold homeomorphic but not diffeomorphic to \( S^n \).  
For a Riemannian metric \( g \) on \( \Sigma \), let \( \mathcal{L}(g) = \{ L(\gamma) : \gamma \text{ closed geodesic of } g \} \subset \mathbb{R}_{>0} \) denote the **length spectrum**, regarded as a set of lengths of distinct closed geodesics.  
Consider the following **asymptotic Zoll condition**:

1. The set of unit tangent vectors tangent to closed geodesics is **topologically dense** in the unit tangent bundle \( U T\Sigma \); that is, for any open subset \( U \subset U T\Sigma \), there exists a closed geodesic whose tangent vector enters \( U \).  
2. The length spectrum \( \mathcal{L}(g) \) possesses a single finite accumulation point \( L_0 > 0 \), meaning there exists a sequence \( L_i \in \mathcal{L}(g) \) with \( L_i \to L_0 \) as \( i \to \infty \) and no other accumulation points.

**Question:**  
Does there exist a Riemannian metric \( g \) on an exotic sphere \( \Sigma^n \) with *positive sectional curvature* that satisfies this asymptotic Zoll condition?  
Equivalently, does the conjunction of positive curvature and “quasi-Zoll” dynamical behavior of the geodesic flow rule out exotic smooth structures on spheres?

Why is it a "good" problem:  
This question explores whether a deep dynamical regularity in geodesic flows can force standard smoothness on a topologically fixed manifold. It extends the spirit of classical differentiable sphere theorems by replacing curvature pinching (a local geometric constraint) with a global dynamical one (density and spectral convergence of closed geodesics).  
The problem directly links three central themes:
- **Positive sectional curvature**, a rare and still poorly understood geometric condition.  
- **Closed geodesic dynamics**, central to Hamiltonian and symplectic geometry.  
- **Exotic smooth structures**, governed by the homotopy sphere groups \( \Theta_n \).  

A solution would illuminate whether dynamical rigidity of geodesic behavior can detect or exclude exotic differentiable structures. The question is simple to state yet reaches the frontier of several difficult subjects—combining geometry, topology, and dynamics into one sharply formulated challenge. This synthesis and potential depth make it a quintessential example of a “good” research problem.